%===============================================================================
% LaTeX-sjabloon voor het rapport van het graduaatsproject
% Graduaat in het Programmeren - HOGENT
%
% Je hoeft in dit bestand normaal niets aan te passen: je gegevens staan in
% ../projectgegevens.tex. Schrijf je tekst in de hoofdstukbestanden.
%===============================================================================

\documentclass[dutch,dit,report]{hogentreport}

% Schrijf je in het Engels? Vervang `dutch' door `english'. Let op: dat kan
% enkel na uitdrukkelijke toestemming van je begeleider, en je Nederlandse
% samenvatting blijft dan verplicht (zie samenvatting.tex).

%% Afbeeldingen komen in de map graphics/
\graphicspath{{../graphics/}}

%% Codefragmenten met kleuraccenten. Compileer met de -shell-escape vlag;
%% de make_report-scripts doen dat al voor jou.
%%
%% Compileer je met een oudere LaTeX-installatie (TeX Live 2022 of ouder, met
%% minted versie 2)? Vervang de regel hieronder dan door:
%%     \usepackage[chapter,outputdir=../output]{minted}
%% Vanaf minted 3 regelt het pakket dat zelf en geeft die optie een fout.
\usepackage[chapter]{minted}
\usemintedstyle{solarized-light}

\setminted{%
    autogobble,
    frame=lines,
    breaklines,
    linenos,
    tabsize=4
}

\renewcommand\listoflistingscaption{%
    \IfLanguageName{dutch}{Lijst van codefragmenten}{List of listings}
}
\renewcommand\listingscaption{%
    \IfLanguageName{dutch}{Codefragment}{Listing}
}
\renewcommand*\listoflistings{%
    \cleardoublepage\phantomsection\addcontentsline{toc}{chapter}{\listoflistingscaption}%
    \listof{listing}{\listoflistingscaption}%
}

%% Klikbare URL's, ook in de tekst
\usepackage{xurl}

%%---------- Je gegevens -------------------------------------------------------
%% Alles wat hieronder ingevuld wordt, komt uit ../projectgegevens.tex
\input{../projectgegevens}

\author{\studentnaam}
\supervisor{\begeleidernaam}
\cosupervisor{\mentornaam}
\title[\projectondertitel]{\projecttitel}
\academicyear{\academiejaarwaarde}
\examperiod{1}          %% 1 = eerste examenperiode, 2 = tweede, 3 = derde
\degreesought{\IfLanguageName{dutch}{\opleidingnaam}{Associate Degree in Programming}}
\institution{\bedrijfsnaam}

%% Afbrekingen die LaTeX verkeerd doet, kan je hier opgeven
\hyphenation{back-slash}

%% Bibliografie
\addbibresource{bronnen.bib}              %% Bronnen van je rapport
\addbibresource{../voorstel/voorstel.bib} %% Bronnen van je projectvoorstel
\defbibheading{bibempty}{}

%% Geen lege bladzijden tussen hoofdstukken
\renewcommand{\cleardoublepage}{\clearpage}

\renewcommand{\arraystretch}{1.2}

\begin{document}

%---------- Voorwerk -----------------------------------------------------------

\frontmatter

\hypersetup{pageanchor=false}
%% Bij een Engelstalig rapport komt er ook een Nederlands titelblad
\IfLanguageName{english}{%
    \degreesought{\opleidingnaam}%
    \begin{otherlanguage}{dutch}%
       \maketitle%
    \end{otherlanguage}%
}{}

\maketitle
\hypersetup{pageanchor=true}

%%=============================================================================
%% Woord vooraf
%%=============================================================================

\chapter*{\IfLanguageName{dutch}{Woord vooraf}{Preface}}%
\label{ch:voorwoord}

%% Richtlengte: een halve bladzijde. Dit is het enige deel van je rapport waar
%% je in de ik-vorm schrijft.
%%
%% Schrijf hier kort:
%% - waarom dit onderwerp jou interesseerde;
%% - wat je ervan verwachtte toen je eraan begon;
%% - wie je wil bedanken (je mentor, je team, je begeleider).
%%
%% Let op: dit is geen samenvatting van je project. Wie enkel dit leest, hoeft
%% nog niet te weten wat je gebouwd hebt.

Vervang deze tekst door je eigen woord vooraf.

%%=============================================================================
%% Samenvatting - VERPLICHT IN TWEE TALEN
%%=============================================================================
%%
%% Je rapport begint met twee samenvattingen, elk op een eigen bladzijde:
%%
%%   bladzijde 1: de Nederlandse samenvatting
%%   bladzijde 2: dezelfde samenvatting in het Engels
%%
%% Beide zijn verplicht, ook als je je rapport volledig in het Nederlands
%% schrijft. Ze staan bovenaan omdat ze in de praktijk het enige zijn dat een
%% druk jurylid met zekerheid leest.
%%
%% Richtlengte: 200 tot 250 woorden per taal, doorlopende tekst, geen lijstjes.
%%
%% Een goede samenvatting beantwoordt in die volgorde:
%%   1. Voor welk bedrijf en welk probleem?
%%   2. Wat heb je gemaakt?
%%   3. Met welke technologie, en waarom die?
%%   4. Wat is het resultaat - werkt het, en hoe weet je dat?
%%   5. Wat heeft het bedrijf eraan?
%%
%% Schrijf ze op het einde, niet in het begin. En schrijf ze zelf: dit is de
%% bladzijde waaraan een lezer merkt of je je eigen werk begrepen hebt.

%%---------- Nederlandse samenvatting -----------------------------------------

\IfLanguageName{english}{\selectlanguage{dutch}}{}

\chapter*{Samenvatting}
\label{ch:samenvatting-nl}

Vervang deze tekst door je Nederlandse samenvatting van ongeveer 200 woorden.

\IfLanguageName{english}{\selectlanguage{english}}{}

%%---------- Engelse samenvatting ---------------------------------------------

\IfLanguageName{dutch}{\selectlanguage{english}}{}

\chapter*{Abstract}
\label{ch:samenvatting-en}

Replace this text with the English version of your summary, about 200 words.
It is a translation of the Dutch summary, not a different text.

\IfLanguageName{dutch}{\selectlanguage{dutch}}{}


%---------- Inhoud, lijsten ----------------------------------------------------

\tableofcontents

% Geef elke figuur, tabel en codefragment een KORT bijschrift tussen vierkante
% haken; enkel dat korte bijschrift komt in de lijsten terecht:
%   \caption[korte omschrijving]{de volledige omschrijving}

\listoffigures

% Heb je geen tabellen of codefragmenten? Zet de bijhorende regel dan in
% commentaar; een lege lijst staat slordig.
\listoftables
\listoflistings

%---------- Kern ---------------------------------------------------------------

\mainmatter{}

% Richtwaarde voor het geheel: 10 tot 15 bladzijden kerntekst. Dat is weinig,
% en dat is met opzet: je rapport is een technisch verslag van een project van
% ongeveer tien werkdagen, geen scriptie. Schrap liever dan te vullen.

TODO!: INLEIDING VOLLEDIG NALEZEN
%%=============================================================================
%% Inleiding
%%=============================================================================

\chapter{\IfLanguageName{dutch}{Inleiding}{Introduction}}%
\label{ch:inleiding}

Vesalius.ai ontwikkelt software die het dagelijkse werk van artsen en andere zorgverleners moet vereenvoudigen. Het platform ondersteunt verschillende processen binnen de medische praktijk en heeft als doel om administratieve handelingen te beperken, zodat zorgverleners zich meer kunnen richten op hun kerntaak. Binnen deze omgeving wordt voortdurend gezocht naar verbeteringen die de configuratie, automatisering en gebruiksvriendelijkheid van het platform verhogen.
Tijdens mijn stage werkte ik mee binnen het ontwikkelingsteam van Vesalius.ai. Mijn stagewerk bestond uit verschillende taken binnen de bestaande applicatie en ontwikkelomgeving. Het graduaatsproject werd daarbinnen afgebakend als een zelfstandig project op basis van bestaande en al langere tijd openstaande JIRA-tickets. Op die manier sluit het graduaatsproject rechtstreeks aan bij een concrete behoefte van het bedrijf, terwijl het tegelijk voldoende ruimte biedt om zelfstandig analyse-, ontwerp-, ontwikkel- en testwerk uit te voeren.
Een bredere aanleiding voor verschillende verbeteringen binnen het platform is de nood aan flexibele en correct gegenereerde communicatie en documenten. Zo staat binnen het team ook al langere tijd het idee op de backlog om e-mails die door het platform worden verstuurd via een WYSIWYG-editor te kunnen beheren. Daarbij vormt de compatibiliteit van gegenereerde HTML met verschillende e-mailclients een belangrijke technische uitdaging. Dit bredere probleem illustreert dat configuratie en automatische generatie van output een belangrijk onderdeel vormen binnen het platform. Het e-mailproject zelf valt echter buiten de scope van dit graduaatsproject.

\section{\IfLanguageName{dutch}{Probleemstelling}{Problem statement}}%
\label{sec:probleemstelling}

Binnen het Vesalius-platform worden documenten en andere vormen van output gedeeltelijk automatisch gegenereerd op basis van geconfigureerde templates. De bestaande werking behandelt een template in bepaalde situaties op organisatieniveau. Hierdoor is het niet mogelijk om voor iedere template aan te geven dat deze uitsluitend relevant is voor één of meerdere specifieke artsen.
In de praktijk kunnen bepaalde documenttemplates echter specifiek bedoeld zijn voor het werkproces van bepaalde artsen. Wanneer dergelijke templates als automatisch te genereren worden ingesteld, kan de huidige werking ertoe leiden dat ook andere artsen output ontvangen die voor hen niet relevant is. Dit veroorzaakt onnodige documenten en vermindert de relevantie van automatisch gegenereerde informatie.
Daarnaast bestaat er voor consultatietemplates van Scribe een behoefte om aan te geven voor welke doelgroep de gegenereerde output bedoeld is. Een patiënt heeft doorgaans nood aan begrijpelijke mensentaal, terwijl een arts eerder medische terminologie verwacht. De gewenste doelgroep moet daarom onderdeel kunnen worden van de configuratie van een consultatietemplate.
De twee problemen zijn vergelijkbaar doordat ze betrekking hebben op het beter afstemmen van automatisch gegenereerde output op de context waarin deze wordt gebruikt.

\section{\IfLanguageName{dutch}{Doelstelling}{Objective}}%
\label{sec:doelstelling}

Het doel van het graduaatsproject is om twee bestaande functionaliteiten binnen het Vesalius-platform uit te breiden, zodat automatisch gegenereerde output gerichter kan worden afgestemd op de gebruiker en de doelgroep.
Voor documenttemplates wordt onderzocht en gerealiseerd hoe een template optioneel aan één of meerdere specifieke artsen kan worden gekoppeld. De automatische generatie moet vervolgens rekening houden met deze koppeling, zonder de algemene zichtbaarheid van de template te wijzigen. Manuele generatie blijft voor alle gebruikers mogelijk.
Voor consultatemplates van Scribe wordt onderzocht en gerealiseerd hoe de gewenste doelgroep van de gegenereerde output kan worden ingesteld. Hierbij wordt onderscheid gemaakt tussen patiëntgerichte output in begrijpelijke taal en artsgerichte output met medische terminologie.
Het einddoel is een oplossing die aansluit bij de bestaande architectuur en ontwikkelstandaarden van Vesalius.ai en die door het team verder onderhouden en uitgebreid kan worden.

\section{\IfLanguageName{dutch}{Afbakening}{Scope}}%
\label{sec:afbakening}

Het graduaatsproject is bewust beperkt tot twee bestaande JIRA-tickets die binnen de werking van het Vesalius.ai-team reeds op de backlog stonden.
Het eerste onderdeel betreft het optioneel koppelen van documenttemplates aan één of meerdere artsen. De selectielijst moet daarbij uitsluitend gebruikers met de functie ``arts'' tonen. De koppeling beïnvloedt enkel het automatische gebruik van een template en niet de zichtbaarheid ervan. Wanneer geen arts gekoppeld is, blijft het bestaande gedrag behouden. Wanneer wel artsen gekoppeld zijn en automatische generatie is ingeschakeld, wordt de template enkel automatisch gegenereerd wanneer een gekoppelde arts is ingelogd. Manuele generatie blijft voor alle gebruikers mogelijk.
Het tweede onderdeel betreft consultatemplates van Scribe. Hierbij wordt per template ingesteld of de gegenereerde output bedoeld is voor een patiënt of voor een arts. Bij patiëntgerichte output ligt de nadruk op begrijpelijke mensentaal; bij artsgerichte output wordt medische terminologie gebruikt.
Scribe- of consultatietemplates worden in het eerste ticket uitdrukkelijk als aparte functionaliteit beschouwd en vallen daar buiten de scope. Voor het tweede ticket vormen de Scribe-consultatemplates net de specifieke scope.
Ook het bredere idee rond een WYSIWYG-editor voor e-mails valt buiten dit graduaatsproject. Dit onderwerp bevindt zich op de bredere productbacklog en vormt geen onderdeel van de twee geselecteerde JIRA-tickets.
De keuzes over de functionele scope en prioriteiten werden binnen het team bepaald op basis van de bestaande JIRA-tickets en de productbehoeften. Binnen deze afgesproken scope heb ik zelfstandig keuzes gemaakt over de technische uitwerking, de structuur van de oplossing, de implementatie en de manier waarop de functionaliteit werd getest.

\section{\IfLanguageName{dutch}{Stagewerk en graduaatsproject}{Internship work and graduation project}}%
\label{sec:stage-vs-project}

Het graduaatsproject maakt deel uit van mijn stage bij Vesalius.ai, maar beide begrippen hebben niet exact dezelfde scope. Tijdens de stage werkte ik mee aan verschillende bestaande taken binnen het ontwikkelingsteam en maakte ik kennis met de ontwikkelomgeving, de codebase, de werkwijze en de bestaande technische infrastructuur van het bedrijf.
Het graduaatsproject werd daarbinnen afgebakend als een zelfstandig technisch project. De functionele scope werd gebaseerd op bestaande JIRA-tickets die reeds langere tijd op de backlog stonden. Hierdoor kon ik werken aan een probleem dat rechtstreeks relevant is voor het bedrijf, terwijl ik tegelijk zelf verantwoordelijk was voor de analyse, ontwerpkeuzes, implementatie en validatie binnen de afgesproken grenzen.
De bestaande bedrijfsarchitectuur, programmeertalen, frameworks, infrastructuur en ontwikkelafspraken waren grotendeels vastgelegd door het team. Mijn zelfstandigheid lag voornamelijk in de analyse van de tickets, het bepalen van de technische aanpak binnen deze omgeving, de concrete implementatie en het controleren van de oplossing.

\section{\IfLanguageName{dutch}{Leeswijzer}{Guide to the report}}%
\label{sec:leeswijzer}

Na deze inleiding wordt in hoofdstuk~\ref{ch:technologie} de technologische context van het project besproken. Daarbij wordt toegelicht welke technologieën reeds binnen Vesalius.ai gebruikt worden en waarom deze technologieën ook voor dit project werden behouden.
Hoofdstuk~\ref{ch:ontwerp} behandelt de analyse en het ontwerp van de oplossing. Hierbij worden de functionele en niet-functionele vereisten, de architectuur, het gegevensmodel, privacy en veiligheid en de belangrijkste ontwerpkeuzes besproken.
In hoofdstuk~\ref{ch:realisatie} wordt de effectieve realisatie beschreven. De codeopbouw, werkwijze, versiebeheer, belangrijkste technische onderdelen en problemen tijdens de ontwikkeling komen daar aan bod.
Hoofdstuk~\ref{ch:resultaat} bespreekt het uiteindelijke resultaat en de manier waarop de functionaliteit werd gevalideerd. Tot slot bevat hoofdstuk~\ref{ch:reflectie} de conclusie, persoonlijke reflectie, geleerde lessen en mogelijke verdere ontwikkelingen.
%%=============================================================================
%% Technologische context en keuze
%%=============================================================================

\chapter{\IfLanguageName{dutch}{Technologie en verantwoording}{Technology and rationale}}%
\label{ch:technologie}

%% Richtlengte: 2 tot 3 bladzijden.
%%
%% LET OP: DIT IS GEEN LITERATUURSTUDIE.
%% Je hoeft geen overzicht te schrijven van wat er in het vakgebied bestaat,
%% en je hoeft geen wetenschappelijke bronnen te verzamelen. De vraag hier is
%% praktisch: welke technologie heb je gebruikt, welke alternatieven waren er,
%% en waarom is het deze geworden?
%%
%% Dit is ook wat de jury je tijdens je verdediging zal vragen. Wie het hier
%% opschrijft, staat daar sterker.
%%
%% WEL: vermeld de bronnen die je echt gebruikt hebt. Voor een project zijn dat
%% meestal officiële documentatie, handleidingen, blogberichten van makers,
%% normen, of een boek. Vermeld ze omdat je lezer ze moet kunnen terugvinden,
%% niet omdat een lijst nu eenmaal lang moet zijn. Zie de uitleg onderaan dit
%% bestand en de voorbeelden in bronnen.bib.

Vervang deze tekst door een korte inleiding op het hoofdstuk: welke
technologische omgeving bestond er al bij het bedrijf, en welke ruimte had je
om te kiezen?

\section{\IfLanguageName{dutch}{Bestaande omgeving}{Existing environment}}%
\label{sec:bestaande-omgeving}

%% Waar stap je in? Welke talen, frameworks, databanken, diensten en
%% werkwijzen waren er al? Wat lag vast en wat mocht je zelf beslissen?

\section{\IfLanguageName{dutch}{Vergelijking van de mogelijkheden}{Comparison of the options}}%
\label{sec:vergelijking}

%% Vergelijk twee tot vier reële mogelijkheden op criteria die er voor dit
%% project toe doen: aansluiting bij de bestaande omgeving, leercurve,
%% ondersteuning, licentie, prestaties, onderhoudbaarheid, kostprijs.
%%
%% Een tabel werkt hier goed. Verzin geen criteria om een tabel te vullen:
%% drie criteria die echt gewogen hebben, zijn beter dan acht die je verzon.

\begin{table}[htbp]
  \centering
  \begin{tabular}{lccc}
    \toprule
    \textbf{Criterium} & \textbf{Optie A} & \textbf{Optie B} & \textbf{Optie C} \\
    \midrule
    Aansluiting bij de bestaande omgeving & ++ & +  & -- \\
    Leercurve binnen tien dagen           & +  & ++ & -- \\
    Ondersteuning en documentatie         & ++ & +  & +  \\
    \bottomrule
  \end{tabular}
  \caption[Vergelijking van de overwogen technologieën]{\label{tab:vergelijking}%
    Vergelijking van de overwogen technologieën. Vervang de criteria en de
    beoordeling door je eigen afweging, en licht ze in de tekst toe.}
\end{table}

\section{\IfLanguageName{dutch}{Vergelijking met de opleiding}{Comparison with the programme}}%
\label{sec:vergelijking-opleiding}

%% Verwacht onderdeel van je verdediging: hoe verhoudt de technologie die je
%% hier gebruikte zich tot wat je in de opleiding gezien hebt? Wat was
%% hetzelfde, wat was nieuw, en wat bleek in de praktijk anders te werken dan
%% in de les?

\section{\IfLanguageName{dutch}{Gemaakte keuze}{The choice made}}%
\label{sec:keuze}

%% Eén heldere paragraaf: dit is het geworden, en dit is de reden. Vermeld
%% ook eerlijk wat je ervoor hebt moeten opgeven.
%%
%% Verwijs naar je bronnen met \autocite{sleutel} of \textcite{sleutel}. De
%% sleutels staan in bronnen.bib. Zet de verwijzing binnen de zin, vlak bij de
%% bewering waar ze bij hoort - niet achteraan de paragraaf. Hieronder een
%% voorbeeld van beide vormen; vervang ze door je eigen bronnen.

Bij het beoordelen van leesbaarheid en onderhoudbaarheid van de code werd
uitgegaan van de gangbare vuistregels uit het vakgebied~\autocite{Martin2008}.
De officiële documentatie van het gekozen framework~\autocite{DockerDocs2026}
diende als voornaamste naslagwerk. \textcite{Fowler2018} gebruik je wanneer de
naam van de auteur zelf een rol speelt in je zin.

%%---------- Bronnen vermelden in een projectrapport ---------------------------
%%
%% Welke bronnen horen in je bronnenlijst?
%%   - officiële documentatie van de talen, frameworks en diensten die je gebruikt
%%   - handleidingen, blogberichten of video's waarop je je aanpak gebaseerd hebt
%%   - normen en standaarden waaraan je je gehouden hebt
%%   - een boek of artikel, als je er een gebruikt hebt
%%   - interne documenten van je bedrijf: vermeld ze, maar neem geen
%%     vertrouwelijke inhoud over zonder toestemming van je mentor
%%
%% Wat er NIET in hoort:
%%   - een gesprek met een AI-hulpmiddel. Dat is geen bron die iemand kan
%%     nagaan; het hoort thuis in de bijlage over je AI-gebruik.
%%   - bronnen die je niet gelezen hebt. Een lange lijst maakt geen indruk;
%%     een lijst met werken die nergens aangehaald worden, wel - de verkeerde.
%%
%% Vermeld bij een webbron ALTIJD de datum waarop je ze geraadpleegd hebt
%% (het veld urldate in bronnen.bib): documentatie verandert.

%%=============================================================================
%% Analyse en ontwerp
%%=============================================================================

\chapter{\IfLanguageName{dutch}{Analyse en ontwerp}{Analysis and design}}%
\label{ch:ontwerp}

%% Richtlengte in TOTALITEIT van deze rubriek: 2 tot 3 bladzijden.
%%
%% Hier toon je dat er nagedacht is voordat er getypt werd. Dit hoofdstuk
%% weegt zwaar: het is waar het verschil zichtbaar wordt tussen een uitgevoerde
%% taak en een zelfstandig project.
TODO: 
Voor de implementatie werd eerst onderzocht welke bestaande functionaliteit door de twee JIRA-tickets wordt uitgebreid en welke gevolgen de nieuwe vereisten hebben voor de frontend, backend en database. Het ontwerp vertrekt daarom vanuit het bestaande gedrag en probeert de nieuwe functionaliteit zo beperkt mogelijk in te bouwen zonder bestaande workflows onnodig te wijzigen.

\section{\IfLanguageName{dutch}{Vereisten}{Requirements}}%
\label{sec:vereisten}
TODO! (lees nog eens na)
%% Wat moest het worden? Onderscheid wat noodzakelijk was van wat wenselijk
%% was, en vermeld bij wie je die vereisten opgehaald hebt (je mentor, het
%% team, de gebruikers).
%%
%% Vermeld ook de niet-functionele vereisten die er echt toe deden:
%% snelheid, veiligheid, privacy, onderhoudbaarheid, beschikbaarheid.
De functionele vereisten werden voornamelijk afgeleid uit de bestaande JIRA-tickets en afgestemd binnen het Vesalius.ai-team.
Voor documenttemplates zijn de belangrijkste vereisten:

\begin{itemize}
    \item Een documenttemplate kan optioneel aan één of meerdere artsen worden gekoppeld.
    \item Enkel gebruikers met de functie ``arts'' kunnen worden geselecteerd.
    \item Een template zonder gekoppelde artsen behoudt het bestaande gedrag.
    \item Een template met gekoppelde artsen wordt bij automatische generatie enkel gebruikt wanneer een gekoppelde arts is ingelogd.
    \item Manuele generatie blijft mogelijk voor alle gebruikers.
    \item De koppeling met artsen verandert de zichtbaarheid van de template niet.
\end{itemize}

Voor Scribe-consultatemplates zijn de belangrijkste vereisten:

\begin{itemize}
    \item Per consultatietemplate moet de doelgroep kunnen worden ingesteld.
    \item De mogelijke doelgroepen zijn patiënt en arts.
    \item Patiëntgerichte output moet begrijpelijk zijn voor een patiënt.
    \item Artsgerichte output mag medische terminologie gebruiken.
\end{itemize}

Naast deze functionele vereisten zijn ook niet-functionele vereisten belangrijk. De uitbreiding moet onderhoudbaar blijven, aansluiten bij de bestaande architectuur en geen onnodige wijzigingen veroorzaken aan bestaande functionaliteit. Omdat het platform met medische informatie werkt, zijn privacy en toegangscontrole eveneens belangrijke aandachtspunten.

\section{\IfLanguageName{dutch}{Architectuur}{Architecture}}%
\label{sec:architectuur}
TODO!
%% Toon de opbouw van je oplossing met een figuur: welke onderdelen zijn er,
%% wat doet elk onderdeel, en hoe praten ze met elkaar? Een goed getekend
%% schema vervangt een bladzijde tekst.
%%
%% Plaats je afbeeldingen in de map graphics/.

De oplossing wordt geïntegreerd in de bestaande architectuur van het Vesalius-platform. De Angular-frontend communiceert met de Laravel-backend. De backend verwerkt de bedrijfslogica en communiceert met MySQL voor persistente gegevens. Authenticatie en gebruikersinformatie worden beheerd via de bestaande Keycloak-integratie.
Een vereenvoudigde voorstelling van de architectuur is:

\begin{figure}[htbp]
  \centering
  \includegraphics[width=0.8\textwidth]{architectuur}
  \caption[Architectuur van de oplossing]{Architectuur van de uitbreiding binnen het Vesalius-platform. De Angular-frontend communiceert met de Laravel-backend, die de bedrijfslogica uitvoert en de benodigde gegevens uit MySQL haalt. Keycloak voorziet de bestaande authenticatie- en gebruikerscontext.}
  \label{fig:architectuur}
\end{figure}

De belangrijkste ontwerpkeuze hierbij is dat de nieuwe functionaliteit geen afzonderlijke applicatie of service vormt. De functionaliteit behoort inhoudelijk tot het bestaande template- en generatieproces en wordt daarom in dezelfde applicatie geïntegreerd.

\section{\IfLanguageName{dutch}{Gegevensmodel}{Data model}}%
\label{sec:gegevensmodel}
TODO!
%% Werk je met gegevens, toon dan hoe ze gestructureerd zijn en waarom.
%% Heeft je project geen noemenswaardig gegevensmodel, verwijder deze sectie
%% dan gerust.
Voor documenttemplates moet een relatie kunnen worden gelegd tussen een template en één of meerdere gebruikers met de functie ``arts''. Omdat één template aan meerdere artsen kan worden gekoppeld en één arts meerdere templates kan gebruiken, gaat het conceptueel om een many-to-many-relatie.
Bij het ontwerpen van deze relatie moet rekening worden gehouden met het bestaande datamodel. De nieuwe koppeling mag de bestaande werking van templates zonder artsselectie niet breken. Daarom moet het ontbreken van gekoppelde artsen een geldige toestand blijven en moet dit geïnterpreteerd worden als het bestaande gedrag.
Voor de Scribe-consultatemplates is een bijkomende configuratiewaarde nodig die de gewenste doelgroep van de output beschrijft. Deze waarde kan conceptueel worden voorgesteld als een beperkte keuze tussen patiëntgerichte en artsgerichte output. Een expliciete beperkte waarde voorkomt dat verschillende vrije tekstwaarden voor hetzelfde concept ontstaan.


\section{\IfLanguageName{dutch}{Privacy, veiligheid en regelgeving}{Privacy, security and regulation}}%
\label{sec:privacy}
TODO (goed nalezen!)
%% Verwijder deze sectie NIET zonder na te denken: bij de meeste projecten is
%% ze van toepassing, en deontologisch werken met oog voor veiligheid en
%% privacy is een leerresultaat van dit opleidingsonderdeel.
%%
%% VERWERKT JE TOEPASSING PERSOONSGEGEVENS?
%% Namen, e-mailadressen, klantnummers, foto's, locaties, logbestanden met
%% IP-adressen: dat zijn persoonsgegevens, en dan geldt de AVG (Algemene
%% Verordening Gegevensbescherming, in het Engels de GDPR). Beschrijf dan kort:
%%   - welke gegevens je verwerkt en waarom je die nodig hebt (dataminimalisatie:
%%     wat je niet nodig hebt, verzamel je niet);
%%   - hoe lang ze bewaard blijven;
%%   - wie ze kan zien, en hoe je dat afschermt;
%%   - of ze het bedrijf verlaten, bijvoorbeeld naar een clouddienst.
%%
%% Werkte je tijdens de ontwikkeling met echte gegevens? Dat hoort in principe
%% niet: gebruik verzonnen of geanonimiseerde testgegevens. Deed je het toch,
%% leg dan uit met welke afspraken.
%%
%% ZIT ER AI IN JE TOEPASSING?
%% Dan is de Europese AI-verordening (de AI Act, verordening 2024/1689) van
%% belang; haar verplichtingen zijn sinds augustus 2026 grotendeels van
%% toepassing. Vermeld in welke risicocategorie je toepassing valt en welke
%% transparantieplichten gelden - een gebruiker moet bijvoorbeeld weten dat hij
%% met een AI-systeem praat of dat inhoud automatisch gegenereerd is.
%% Zit er geen AI in je toepassing, dan volstaat één zin die dat zegt.
%%
%% VEILIGHEID
%% Wachtwoorden, sleutels en verbindingsgegevens horen niet in je repository,
%% en niet in de schermafbeeldingen in dit rapport. Beschrijf kort hoe je
%% toegang en geheimen geregeld hebt.
Het Vesalius-platform verwerkt medische en andere persoonsgegevens. Daarom moet de uitbreiding rekening houden met de principes van gegevensbescherming en toegangscontrole. De nieuwe functionaliteit moet uitsluitend gebruikmaken van de gebruikersgegevens die noodzakelijk zijn om de templateconfiguratie en automatische generatie correct uit te voeren.
Voor de artsselectie is het bijvoorbeeld niet nodig om bijkomende persoonsgegevens op te slaan wanneer de bestaande gebruikersidentiteit reeds beschikbaar is. De koppeling kan gebruikmaken van de bestaande gebruikersidentificatie binnen het platform.
Authenticatie wordt niet opnieuw geïmplementeerd binnen dit project. Keycloak blijft verantwoordelijk voor de bestaande authenticatie- en identiteitscontext. De applicatie moet vervolgens de bestaande autorisatie- en gebruikersgegevens correct gebruiken.
Tijdens ontwikkeling en testen moet bovendien vermeden worden dat echte medische patiëntgegevens onnodig worden gebruikt. Testscenario's kunnen uitgevoerd worden met test- of geanonimiseerde gegevens.
Het project bevat geen nieuwe autonome AI-functionaliteit. De configuratie van patiënt- of artsgerichte output bepaalt wel hoe bestaande gegenereerde inhoud moet worden afgestemd, maar het graduaatsproject introduceert geen afzonderlijk AI-model.
Geheimen, wachtwoorden, tokens en andere gevoelige configuratiegegevens worden niet opgenomen in de broncode of screenshots van het rapport.

\section{\IfLanguageName{dutch}{Ontwerpkeuzes}{Design decisions}}%
\label{sec:ontwerpkeuzes}
TODO!+rewrite 
%% Noteer drie tot vijf keuzes die je bewust gemaakt hebt, telkens in deze
%% vorm: de keuze, de alternatieven, de reden, en het gevolg.
%%
%% Dit is de sectie waaruit een jurylid zijn vragen haalt. Een keuze waarvan
%% je het alternatief niet kan benoemen, was geen keuze.
1) om de artskoppeling optioneel te maken. Een alternatief zou zijn om voor elke template verplicht één of meerdere artsen te selecteren. Dat zou echter het bestaande gedrag wijzigen voor templates die vandaag organisatiebreed worden gebruikt. Door de koppeling optioneel te maken, blijft backwards compatibility behouden.
2) om de artskoppeling enkel invloed te laten hebben op automatische generatie en niet op zichtbaarheid. Een alternatief zou zijn om templates die aan bepaalde artsen gekoppeld zijn ook voor andere gebruikers te verbergen. Dat zou echter twee verschillende concepten door elkaar halen: wie een template mag zien en voor wie een template automatisch wordt gebruikt. Door deze verantwoordelijkheden gescheiden te houden, blijft het gedrag duidelijker.
3) om enkel gebruikers met de functie ``arts'' te tonen in de selectie. Een vrije gebruikersselectie zou technisch eenvoudiger kunnen lijken, maar zou functioneel incorrecte configuraties mogelijk maken. De selectie wordt daarom beperkt op basis van de bestaande gebruikersfunctie.
4) om voor Scribe een expliciete doelgroepinstelling te gebruiken in plaats van vrije tekst. De alternatieve aanpak waarbij een gebruiker zelf instructies of een doelgroep beschrijft, zou meer flexibiliteit bieden maar ook meer variatie en onvoorspelbaarheid veroorzaken. Een beperkte keuze tussen patiënt en arts maakt de configuratie duidelijker en eenvoudiger te valideren.
5) om de nieuwe functionaliteit binnen de bestaande Laravel- en Angulararchitectuur te implementeren in plaats van een afzonderlijke service te bouwen. Hierdoor blijven gegevens, authenticatie, deployment en onderhoud binnen dezelfde technische context.
%%=============================================================================
%% Realisatie
%%=============================================================================
\chapter{\IfLanguageName{dutch}{Realisatie}{Implementation}}%
\label{ch:realisatie}
Dit hoofdstuk beschrijft hoe de functionele en technische vereisten van het
graduaatsproject worden vertaald naar een zelfstandige proof of concept (PoC)
van een visuele e-maileditor. De nadruk ligt op de belangrijkste onderdelen
van de oplossing, de gemaakte technische keuzes en de aandachtspunten bij de
verwerking van e-mailtemplates.

De realisatie wordt besproken aan de hand van de codeopbouw, de werkwijze en
het versiebeheer, de belangrijkste technische onderdelen en de problemen die
tijdens de ontwikkeling worden onderzocht. De concrete implementatie en de
uiteindelijke resultaten worden beschreven op basis van de effectief
gerealiseerde functionaliteiten.

\section{\IfLanguageName{dutch}{Opbouw van de code}{Structure of the code}}%
\label{sec:opbouw-code}
De proof of concept wordt opgezet als een zelfstandige toepassing, los van de
productieomgeving van Vesalius.ai. Hierdoor kunnen de editor, de verwerking
van e-mailtemplates en de gegenereerde uitvoer afzonderlijk worden ontwikkeld
en getest, zonder de bestaande e-mailverwerking in productie te wijzigen.

De beoogde architectuur bestaat uit een frontend, een mock-API en een
verwerkingslaag voor e-mailtemplates. De frontend wordt ontwikkeld met
Angular en TypeScript. De mock-API wordt opgezet met Node.js, TypeScript en
Express en staat in voor de communicatie tussen de frontend en de logica voor
templatebeheer.

Binnen de frontend bevindt zich de gebruikersinterface waarmee een gebruiker
een e-mailtemplate kan selecteren, de inhoud visueel kan aanpassen en het
resultaat vooraf kan bekijken. Voor de WYSIWYG-functionaliteit wordt Quill 2
als kandidaat onderzocht. De definitieve keuze wordt bepaald door de
technische haalbaarheid en de ondersteuning van de vereiste
bewerkingsmogelijkheden.

De backendlaag behandelt het laden en opslaan van templates en coördineert
de verwerking van de inhoud. De opslag wordt zo opgezet dat de PoC niet
onnodig afhankelijk wordt van één specifieke databasetechnologie. Afhankelijk
van de implementatie kunnen hiervoor tijdelijke opslag of vooraf gedefinieerde
testgegevens worden gebruikt.

De renderingpipeline vormt een afzonderlijk onderdeel van de toepassing.
Deze laag verwerkt de HTML-inhoud, controleert toegestane placeholders en
combineert de bewerkbare inhoud met de vaste structuur van de e-mail. Indien
nodig wordt de CSS inline geplaatst om de compatibiliteit met e-mailclients
te verbeteren.

De lokale testomgeving ondersteunt het controleren van de gegenereerde
e-mails. Hiervoor wordt Mailpit als mogelijke testmailservice gebruikt.
Docker kan worden ingezet om de benodigde services reproduceerbaar op te
starten.

Door de frontend, de API, de templateopslag en de renderingpipeline logisch
van elkaar te scheiden, blijven de verantwoordelijkheden duidelijk. Dit
vereenvoudigt het afzonderlijk testen van de onderdelen en maakt het
eenvoudiger om de PoC later technisch te evalueren met het oog op een
mogelijke integratie in Vesalius.ai.

\section{\IfLanguageName{dutch}{Werkwijze en versiebeheer}{Way of working and version control}}%
\label{sec:werkwijze}
De werkzaamheden worden opgevolgd via Jira. De beschikbare tickets en
projectvereisten vormen het uitgangspunt voor de planning en de uitvoering
van de werkzaamheden. Per onderdeel wordt bepaald welke functionaliteit nodig
is, welke technische aandachtspunten onderzocht moeten worden en hoe het
resultaat kan worden gecontroleerd.

De broncode wordt beheerd met Git en Bitbucket binnen de interne
ontwikkelomgeving van Vesalius.ai. Hierdoor kunnen wijzigingen worden
opgevolgd en kan de evolutie van de implementatie worden nagegaan. De
repository is intern en wordt niet als publieke broncode beschikbaar gesteld.

Voor de ontwikkeling wordt Visual Studio Code gebruikt. De lokale
ontwikkelomgeving draait op Windows 11 met Ubuntu via WSL2. Docker wordt
gebruikt binnen de ontwikkelworkflow wanneer containers nodig zijn voor de
testomgeving of ondersteunende services.

De ontwikkeling volgt een iteratieve werkwijze. Eerst worden de vereisten en
de technische haalbaarheid onderzocht. Vervolgens wordt de basisstructuur van
de toepassing opgezet, waarna de editor, het templatebeheer en de
renderingpipeline stapsgewijs worden uitgewerkt. Daarna worden de
previewfunctionaliteit, de testmailomgeving en de geautomatiseerde tests
toegevoegd volgens de projectplanning.

Na iedere relevante wijziging wordt gecontroleerd of de aangepaste
functionaliteit aan de verwachte vereisten voldoet. Daarbij wordt niet alleen
rekening gehouden met de zichtbare werking van de editor, maar ook met de
HTML-uitvoer, de verwerking van placeholders en de veiligheid van de
gegenereerde e-mail.

De precieze voortgang, de uitgevoerde taken en eventuele afwijkingen worden
bijgehouden aan de hand van de projectplanning en de bijbehorende tickets.
De definitieve werkwijze en de werkelijk uitgevoerde controles worden in dit
hoofdstuk aangevuld op basis van de gerealiseerde implementatie.

\section{\IfLanguageName{dutch}{Belangrijkste onderdelen}{Key components}}%
\label{sec:onderdelen}
De realisatie is opgebouwd rond drie belangrijke technische onderdelen:
de visuele editor en het templatebeheer, de renderingpipeline en de
testomgeving voor de gegenereerde e-mails.

\subsection{Visuele editor en templatebeheer}
De visuele editor vormt het centrale onderdeel van de PoC. Het doel is om
gebruikers toe te laten de inhoud en ondersteunde opmaak van vooraf bepaalde
e-mailtemplates aan te passen zonder rechtstreeks HTML-code te moeten
bewerken.

De gebruiker selecteert een van de drie piloottypes: een
afspraakbevestiging, een afspraakherinnering of een afspraakannulering.
Vervolgens wordt de bijbehorende template geladen in de editor. De gebruiker
kan de beschikbare inhoud aanpassen en het resultaat vooraf bekijken.

De editor richt zich op de inhoud van de e-mail. De vaste e-mailschil, met
bijvoorbeeld het logo, de merkkleuren en de footer, blijft grotendeels
buiten de bewerkingsmogelijkheden. Hierdoor blijft de scope van de PoC
beperkt en kan de aandacht uitgaan naar het aanpassen en betrouwbaar
renderen van de e-mailinhoud.

Naast het bewerken is ook het opslaan en opnieuw laden van templates van
belang. De opslaglaag wordt daarom gescheiden van de editorcomponent.
Hierdoor hoeft de interface niet rechtstreeks afhankelijk te zijn van de
gekozen opslagtechnologie.

De concrete implementatie wordt geëvalueerd op basis van het laden van
templates, het aanpassen van inhoud, het bewaren van wijzigingen en het
opnieuw openen van een opgeslagen template.

\subsection{Renderingpipeline en veilige verwerking}
Een tweede belangrijk onderdeel is de omzetting van de bewerkte template
naar uiteindelijke e-mailinhoud. De HTML die door een visuele editor wordt
gegenereerd, kan niet zonder controle worden gebruikt. De invoer moet worden
gevalideerd en niet-toegestane HTML-elementen of attributen moeten worden
verwijderd of geweigerd.

De renderingpipeline verwerkt daarom de inhoud volgens vooraf bepaalde
regels. Eerst wordt de HTML gesanitiseerd. Vervolgens worden uitsluitend
toegestane placeholders verwerkt met fictieve testgegevens. Daarna wordt de
inhoud gecombineerd met de vaste e-mailschil. Indien nodig wordt CSS inline
geplaatst om de weergave in e-mailclients te ondersteunen.

Bij de placeholderverwerking is het belangrijk dat onbekende of ontbrekende
waarden veilig worden afgehandeld. Ook moet rekening worden gehouden met de
context waarin een waarde wordt ingevoegd. Het ontsnappen van waarden en het
controleren van toegestane invoer blijven noodzakelijk om ongewenste HTML-
of scriptinhoud te vermijden.

De renderingpipeline wordt afzonderlijk beoordeeld van de editor. Een
correct werkende editor garandeert immers niet dat de uiteindelijke HTML
veilig is of in iedere e-mailclient hetzelfde wordt weergegeven.

\subsection{Preview, testmails en compatibiliteitscontroles}
Het derde onderdeel omvat de controle van de gegenereerde e-mail. De
previewfunctie maakt het mogelijk om de gerenderde inhoud te beoordelen
voordat een testmail wordt verstuurd. Daarbij wordt nagegaan of de inhoud,
de placeholders en de vaste e-mailschil correct worden weergegeven.

Voor het testen van e-mails wordt een lokale testomgeving voorzien. Mailpit
kan worden gebruikt om testmails op te vangen en te inspecteren zonder dat
deze naar echte ontvangers worden verstuurd. De testgegevens zijn fictief
en bevatten geen echte patiëntinformatie.

Daarnaast worden geautomatiseerde tests voorzien voor onderdelen zoals
templatevalidatie, placeholderverwerking en HTML-sanitatie. Vitest kan
worden ingezet voor unit- en integratietests en Playwright voor
browsergebaseerde tests, afhankelijk van de uiteindelijke implementatie.

Ten slotte wordt de weergave van de gegenereerde e-mails onderzocht in
beschikbare e-mailclients, waaronder Gmail, Outlook en Apple Mail. Eventuele
verschillen in HTML- en CSS-ondersteuning worden geregistreerd. De resultaten
van deze controles worden opgenomen in de bijlage met testresultaten.

\subsection{Broncode en vertrouwelijkheid}
De broncode van de toepassing bevindt zich in de interne Bitbucket-omgeving
van Vesalius.ai. Omdat de repository deel uitmaakt van de bedrijfsomgeving,
worden interne bronbestanden, bedrijfsgegevens en vertrouwelijke configuratie
niet zonder toestemming opgenomen in dit rapport.

De technische werking kan worden toegelicht met architectuurschema's,
beschrijvingen van de gegevensstroom en vereenvoudigde codefragmenten.
Wanneer een codefragment uit de interne implementatie niet mag worden
gepubliceerd, wordt de relevante logica op een hoger niveau beschreven.
Zo blijft de technische onderbouwing van het project mogelijk zonder
vertrouwelijke bedrijfsinformatie prijs te geven.

\section{\IfLanguageName{dutch}{Problemen onderweg}{Problems along the way}}%
\label{sec:problemen}
Een belangrijk aandachtspunt bij de realisatie is de betrouwbaarheid van de
gegenereerde e-mail. Een WYSIWYG-editor kan inhoud visueel correct weergeven,
maar de resulterende HTML moet ook veilig worden verwerkt en bruikbaar zijn
in verschillende e-mailclients. Dit vraagt om een afzonderlijke controle van
de editoruitvoer en de uiteindelijke rendering.

Een tweede aandachtspunt is het correct verwerken van placeholders. Dynamische
waarden moeten gecontroleerd worden ingevuld en mogen niet leiden tot
onverwachte HTML-uitvoer. Daarom moet vooraf worden bepaald welke placeholders
toegestaan zijn en hoe ontbrekende of onbekende waarden worden afgehandeld.

Daarnaast is de ondersteuning van HTML en CSS niet identiek in alle
e-mailclients. Een e-mail die in de browserpreview correct wordt weergegeven,
kan in Outlook, Gmail of Apple Mail afwijkingen vertonen. De preview alleen
volstaat daarom niet als bewijs van compatibiliteit. Waar mogelijk worden
gegenereerde testmails ook in de beschikbare e-mailclients gecontroleerd.

Ook de keuze van de opslagmethode vraagt aandacht. Omdat het project een
zelfstandige PoC is en geen productie-integratie, moet de opslag voldoen aan
de testvereisten zonder onnodige afhankelijkheden te introduceren. De
opslaglaag wordt daarom zo veel mogelijk losgekoppeld van de rest van de
toepassing.

Tot slot moet de scope van het project bewaakt worden. De PoC is bedoeld om
de haalbaarheid van een visuele e-maileditor aan te tonen en niet om de
bestaande e-mailinfrastructuur van Vesalius.ai volledig te vervangen.
Productieverzending, integratie met de bestaande platformarchitectuur en
verwerking van echte patiëntgegevens vallen buiten de scope.

De uiteindelijke bespreking van problemen wordt aangevuld met de concrete
technische obstakels die tijdens de implementatie werden vastgesteld, de
oplossingen die daarvoor werden gekozen en de resultaten van de uitgevoerde
hertests.

%%=============================================================================
%% Resultaat en validatie
%%=============================================================================

%%TODO: toevoegen screenshots/testresults/conclusie result

\chapter{\IfLanguageName{dutch}{Resultaat en validatie}{Result and validation}}%
\label{ch:resultaat}
Dit hoofdstuk beschrijft het beoogde resultaat van het graduaatsproject en
de manier waarop de gerealiseerde proof of concept (PoC) wordt gevalideerd.
Het project richt zich op de ontwikkeling van een zelfstandige visuele
e-maileditor waarmee gebruikers de inhoud van bepaalde automatisch
gegenereerde e-mails kunnen aanpassen, opslaan, vooraf bekijken en testen.

Het beoogde resultaat is een werkende demonstratie waarin de verschillende
onderdelen van dit proces samenkomen. De gebruiker moet een bestaande
e-mailtemplate kunnen selecteren, de inhoud visueel kunnen aanpassen en
het resultaat kunnen beoordelen voordat een testmail wordt verstuurd.
Daarnaast moet de oplossing aantonen dat de gegenereerde HTML gecontroleerd
wordt verwerkt en geschikt kan worden gemaakt voor weergave in gangbare
e-mailclients.

De validatie richt zich zowel op de functionele werking als op de technische
betrouwbaarheid van de oplossing. Hierbij worden de vooraf bepaalde vereisten
vergeleken met de werkelijk gerealiseerde functionaliteiten en de resultaten
van de uitgevoerde tests.

\section{\IfLanguageName{dutch}{Wat er opgeleverd is}{What was delivered}}%
\label{sec:opgeleverd}
Het beoogde eindresultaat is een zelfstandige proof of concept van een
WYSIWYG-e-maileditor. De toepassing wordt ontwikkeld buiten de
productieomgeving van Vesalius.ai en dient om de technische haalbaarheid
van visueel aanpasbare e-mailtemplates te onderzoeken.

De PoC richt zich op drie piloottypes: afspraakbevestigingen,
afspraakherinneringen en afspraakannuleringen. Voor deze e-mails moet de
gebruiker de beschikbare inhoud en ondersteunde opmaak kunnen aanpassen
zonder rechtstreeks HTML-code te hoeven bewerken.

Naast de editor omvat het beoogde resultaat een mechanisme om templates
te laden en op te slaan. De wijzigingen moeten opnieuw beschikbaar zijn
wanneer de gebruiker een opgeslagen template opnieuw opent, binnen de
beperkingen van de gekozen opslagmethode.

Een tweede onderdeel is de renderingpipeline. Deze verwerkt de bewerkte
inhoud, controleert de HTML en verwerkt toegestane placeholders met
fictieve gegevens. Vervolgens wordt de inhoud gecombineerd met de vaste
e-mailschil, waarin onder meer het logo, de merkkleuren en de footer kunnen
zijn opgenomen. Indien nodig wordt CSS inline geplaatst om de weergave in
e-mailclients te ondersteunen.

De PoC moet daarnaast een previewfunctie aanbieden waarmee de gebruiker de
gerenderde e-mail kan controleren. Voor verdere validatie wordt een lokale
testmailomgeving voorzien, zodat de gegenereerde e-mails kunnen worden
verstuurd en geïnspecteerd zonder echte ontvangers of patiëntgegevens te
gebruiken.

Het beoogde eindresultaat bestaat bijgevolg uit de volgende onderdelen:

\begin{itemize}
    \item een visuele editor voor de inhoud en ondersteunde opmaak van
    e-mailtemplates;
    \item functionaliteit om templates te laden, aan te passen en op te slaan;
    \item een renderingpipeline voor HTML, placeholders en de vaste e-mailschil;
    \item een previewfunctie voor de gegenereerde e-mail;
    \item een lokale testmailomgeving voor het controleren van de uitvoer;
    \item geautomatiseerde tests en een evaluatie van de compatibiliteit
    met beschikbare e-mailclients;
    \item technische documentatie over de architectuur, de gemaakte keuzes
    en de beperkingen van de PoC.
\end{itemize}

In de definitieve versie van dit hoofdstuk worden screenshots opgenomen
die de werkelijk gerealiseerde functionaliteiten aantonen. Bij voorkeur
wordt de editor samen met een aangepaste template getoond, gevolgd door
een voorbeeld van de gerenderde e-mail. Indien beschikbaar kan ook een
screenshot van de ontvangen testmail in de lokale testomgeving worden
toegevoegd.

\section{\IfLanguageName{dutch}{Testen}{Testing}}%
\label{sec:testen}
De validatie moet aantonen in welke mate de PoC voldoet aan de vooraf
bepaalde functionele en technische vereisten. Hiervoor worden concrete
testscenario's opgesteld die aansluiten bij de belangrijkste gebruikersflows
en onderdelen van de renderingpipeline.

De functionele tests controleren of de gebruiker een template kan selecteren,
de inhoud kan aanpassen, wijzigingen kan opslaan en het resultaat opnieuw
kan laden. Daarnaast wordt nagegaan of de preview de verwachte inhoud
weergeeft en of de testmailfunctionaliteit correct samenwerkt met de lokale
testomgeving.

Voor de templateverwerking zijn onder meer de volgende scenario's relevant:

\begin{itemize}
    \item Een beschikbare e-mailtemplate kan worden geselecteerd en geladen.
    \item De gebruiker kan de inhoud en ondersteunde opmaak aanpassen.
    \item Wijzigingen kunnen worden opgeslagen en opnieuw worden geladen.
    \item De preview toont de verwachte gerenderde inhoud.
    \item Geldige placeholders worden vervangen door de overeenkomstige
    fictieve waarden.
    \item Onbekende of ontbrekende placeholders worden veilig afgehandeld.
    \item De vaste e-mailschil wordt correct gecombineerd met de bewerkbare
    inhoud.
    \item Een testmail kan via de lokale testomgeving worden gecontroleerd.
\end{itemize}

Naast de functionele werking wordt de veiligheid van de HTML-verwerking
gecontroleerd. Hierbij wordt nagegaan of niet-toegestane HTML-elementen
en attributen worden verwijderd of geweigerd en of dynamische waarden
veilig worden verwerkt. Ook wordt onderzocht of de renderingpipeline
onverwachte of onvolledige invoer correct afhandelt.

De compatibiliteitstests richten zich op de weergave van de gegenereerde
e-mails in Gmail, Outlook en Apple Mail, voor zover deze clients beschikbaar
zijn in de testomgeving. Hierbij worden onder meer de inhoud, de opmaak,
de vaste e-mailschil en eventuele afwijkingen gecontroleerd. Wanneer een
client niet beschikbaar is, wordt de compatibiliteit voor die client als
niet-gevalideerd vermeld.

Waar mogelijk worden de tests geautomatiseerd. Unit- en integratietests
kunnen worden gebruikt om afzonderlijke onderdelen van de verwerking te
controleren. Browsertests kunnen daarnaast de belangrijkste gebruikersflows
evalueren.

De concrete testresultaten worden geregistreerd in een afzonderlijke
testtabel. Per scenario worden de verwachte uitkomst, de werkelijke
uitkomst en de status vermeld. Hierdoor kan worden aangetoond welke
vereisten zijn geslaagd, welke problemen werden vastgesteld en welke
onderdelen bijkomende aanpassingen vereisen. De gedetailleerde resultaten
worden opgenomen in Bijlage~\ref{app:testresultaten}.

\section{\IfLanguageName{dutch}{Terugkoppeling van het bedrijf}{Feedback from the company}}%
\label{sec:terugkoppeling}
De terugkoppeling van Vesalius.ai is relevant om te beoordelen of de
oplossing aansluit bij de behoeften van het bedrijf en of de gekozen
technische aanpak bruikbaar is als basis voor verdere ontwikkeling.

Tijdens de evaluatie kan worden nagegaan of de editor voldoende
bewerkingsmogelijkheden aanbiedt, of de afbakening tussen de bewerkbare
inhoud en de vaste e-mailschil duidelijk is en of de gegenereerde e-mails
voldoen aan de verwachte functionele vereisten.

Ook de onderhoudbaarheid van de oplossing is van belang. De scheiding tussen
de frontend, de mock-API, de opslag en de renderingpipeline moet het mogelijk
maken om de afzonderlijke onderdelen te begrijpen en verder te ontwikkelen.
Daarnaast moet worden beoordeeld of de gemaakte keuzes aansluiten bij de
bestaande technische context van Vesalius.ai.

De concrete feedback van de bedrijfsmentor en de betrokken teamleden wordt
in de definitieve versie van dit hoofdstuk opgenomen zodra deze beschikbaar
is. Daarbij worden de belangrijkste opmerkingen, eventuele voorgestelde
aanpassingen en de mogelijke vervolgstappen beschreven.

De PoC is in de eerste plaats bedoeld als technische demonstratie. Een
positieve evaluatie betekent daarom niet automatisch dat de oplossing
klaar is voor productiegebruik. Daarvoor is een afzonderlijke beoordeling
van de integratie, beveiliging, opslag, het beheer en de operationele
vereisten noodzakelijk.

\section{\IfLanguageName{dutch}{Beperkingen}{Limitations}}%
\label{sec:beperkingen}
De belangrijkste beperking van het project is de afgebakende scope. De PoC
richt zich op drie specifieke e-mailtypes en op het aanpassen van de
e-mailinhoud. Het volledig vrij ontwerpen van e-maillay-outs valt buiten
de doelstelling. De vaste e-mailschil blijft grotendeels behouden.

Daarnaast wordt de toepassing niet geïntegreerd in de productieomgeving
van Vesalius.ai. De mock-API en de gekozen opslagmethode dienen om de
functionaliteiten in een gecontroleerde testomgeving te demonstreren.
Ze vormen niet noodzakelijk een geschikte oplossing voor productiegebruik.

Ook de testgegevens zijn bewust beperkt tot fictieve informatie. Hierdoor
kunnen de belangrijkste verwerkingsstappen worden onderzocht zonder echte
patiëntgegevens te gebruiken. De PoC toont daarmee niet automatisch aan
dat alle vereisten voor een productieomgeving zijn ingevuld.

Een verdere beperking betreft de compatibiliteit met e-mailclients.
HTML en CSS kunnen verschillend worden geïnterpreteerd door Gmail, Outlook
en Apple Mail. Zelfs wanneer de preview en de uitgevoerde tests correct
verlopen, kan niet zonder bijkomende validatie worden gegarandeerd dat
iedere e-mailclient de inhoud identiek weergeeft.

Ten slotte hangt de geldigheid van de conclusies af van de werkelijk
uitgevoerde tests en de beschikbare testomgeving. Niet-uitgevoerde
scenario's en niet-beschikbare e-mailclients worden daarom expliciet als
beperkingen vermeld.

De resultaten moeten bijgevolg worden geïnterpreteerd als een evaluatie van
de gerealiseerde PoC binnen de onderzochte context. Ze vormen een basis
om te bepalen welke onderdelen voldoende zijn aangetoond en welke verdere
ontwikkeling of validatie nodig hebben voordat een mogelijke integratie
in Vesalius.ai kan worden overwogen.

%%=============================================================================
%% Conclusie en reflectie
%%=============================================================================

\chapter{\IfLanguageName{dutch}{Conclusie en reflectie}{Conclusion and reflection}}%
\label{ch:reflectie}
Het graduaatsproject richt zich op de ontwikkeling van een zelfstandige proof
of concept (PoC) van een visuele e-maileditor voor Vesalius.ai. Het project
combineert het analyseren van een technisch probleem met het onderzoeken,
ontwerpen en realiseren van een oplossing binnen een professionele
ontwikkelcontext.

De focus ligt op het visueel aanpassen van automatisch gegenereerde
e-mailtemplates, het veilig verwerken van de inhoud en het controleren van
de gegenereerde e-mails in gangbare e-mailclients. Daarbij wordt rekening
gehouden met de beperkingen van e-mailtechnologieën en met de noodzaak om
de oplossing gescheiden te houden van de productieomgeving.

\section{\IfLanguageName{dutch}{Conclusie}{Conclusion}}%
\label{sec:conclusie}
De doelstelling van het graduaatsproject is het ontwikkelen van een
zelfstandige PoC waarmee gebruikers de inhoud en ondersteunde opmaak van
bepaalde e-mailtemplates visueel kunnen aanpassen, opslaan, vooraf bekijken
en testen. Daarnaast wordt onderzocht hoe de gegenereerde HTML veilig kan
worden verwerkt en hoe de e-mails worden weergegeven in verschillende
e-mailclients.

De PoC richt zich op drie piloottypes: afspraakbevestigingen,
afspraakherinneringen en afspraakannuleringen. De vaste e-mailschil, met
bijvoorbeeld het logo, de merkkleuren en de footer, blijft grotendeels
behouden. De nadruk ligt op de bewerkbare inhoud en op een betrouwbare
verwerking van de uiteindelijke e-mail.

De meerwaarde voor Vesalius.ai bestaat erin dat het project de technische
haalbaarheid van visueel aanpasbare e-mailtemplates onderzoekt. Een
geslaagde PoC kan aantonen welke onderdelen nodig zijn om gebruikers meer
controle te geven over de inhoud van e-mails, zonder dat hiervoor onmiddellijk
een integratie in het bestaande productieplatform nodig is.

De uiteindelijke beoordeling van de doelstelling gebeurt aan de hand van
de gerealiseerde functionaliteiten, de geautomatiseerde tests en de
compatibiliteitscontroles. Op basis daarvan kan worden vastgesteld welke
vereisten volledig zijn gerealiseerd, welke slechts gedeeltelijk zijn
aangetoond en welke verdere ontwikkeling vereisen. De PoC vormt daarmee
een mogelijke technische basis voor een toekomstige integratie, maar is
zelf geen productieklare oplossing.

\section{\IfLanguageName{dutch}{Wat ik geleerd heb}{What I learned}}%
\label{sec:geleerd}
Een belangrijk leerpunt is het belang van een duidelijke afbakening bij
softwareontwikkeling. Een visuele e-maileditor lijkt op het eerste gezicht
een afgebakende gebruikersfunctionaliteit, maar omvat ook templatebeheer,
HTML-verwerking, dynamische gegevens, CSS en compatibiliteit met
verschillende e-mailclients. Door deze onderdelen afzonderlijk te
onderzoeken, kan de complexiteit beter worden ingeschat en kan de oplossing
stapsgewijs worden opgebouwd.

Daarnaast leer ik dat een technische keuze niet alleen op basis van
gebruiksgemak moet worden gemaakt. Een editor moet bijvoorbeeld niet enkel
voldoende bewerkingsmogelijkheden aanbieden, maar ook HTML genereren die
veilig kan worden verwerkt en geschikt is voor gebruik in e-mails. De
geschiktheid van een technologie moet daarom worden beoordeeld aan de hand
van concrete vereisten en technische tests.

Ook het onderscheid tussen een preview en de uiteindelijke weergave van
een e-mail is belangrijk. Een e-mail die correct wordt weergegeven in een
browser, wordt niet noodzakelijk identiek weergegeven in Gmail, Outlook of
Apple Mail. Dit maakt duidelijk waarom compatibiliteitscontroles een
afzonderlijk onderdeel van het ontwikkelproces moeten vormen.

Op professioneel vlak biedt het project de mogelijkheid om ervaring op te
doen met het opdelen van een probleem in afzonderlijke technische onderdelen,
het documenteren van ontwerpkeuzes en het systematisch controleren van
vereisten. De werkzaamheden worden opgevolgd via Jira en de broncode wordt
beheerd met Git en Bitbucket. De bestaande bedrijfscontext vraagt bovendien
aandacht voor vertrouwelijkheid en voor het correct omgaan met interne
broncode en configuratiegegevens.

De verdere reflectie wordt aangevuld op basis van de werkelijk uitgevoerde
werkzaamheden, de technische problemen die tijdens de ontwikkeling werden
vastgesteld en de resultaten van de uitgevoerde tests.

\section{\IfLanguageName{dutch}{Wat ik anders zou doen}{What I would do differently}}%
\label{sec:anders}
Bij een gelijkaardig project zou ik de technische risico's zo vroeg mogelijk
onderzoeken. Voor een WYSIWYG-e-maileditor betekent dit onder meer dat ik
vroegtijdig zou testen welke HTML de gekozen editor genereert, welke
bewerkingsmogelijkheden nodig zijn en welke beperkingen zich voordoen bij
de weergave in verschillende e-mailclients.

Daarnaast zou ik de functionele vereisten vanaf het begin vertalen naar
concrete testscenario's. Voorbeelden hiervan zijn het aanpassen en opnieuw
laden van een template, het veilig verwerken van placeholders en het
controleren van de gerenderde HTML. Hierdoor kan tijdens de ontwikkeling
sneller worden vastgesteld of een onderdeel aan de verwachtingen voldoet.

Ook zou ik ontwerpbeslissingen, technische risico's en openstaande vragen
systematisch documenteren. Dit maakt het eenvoudiger om achteraf te
verklaren waarom een bepaalde technologie of architectuur werd gekozen.
Het helpt bovendien om de voortgang te evalueren en om te vermijden dat
belangrijke beslissingen uitsluitend in gesprekken of persoonlijke notities
worden bijgehouden.

Ten slotte zou ik voldoende tijd reserveren voor testen in echte
e-mailclients. De implementatie van de editor en de renderingpipeline is
slechts een deel van het werk; de uiteindelijke bruikbaarheid hangt ook af
van de manier waarop de gegenereerde e-mails worden weergegeven.

\section{\IfLanguageName{dutch}{Verder werk}{Future work}}%
\label{sec:verder-werk}
Een mogelijke volgende stap is de integratie van de PoC in de bestaande
architectuur van Vesalius.ai. Hiervoor moet eerst worden onderzocht hoe de
oplossing aansluit op het bestaande templatebeheer, de e-mailverwerking,
de opslag en de geldende beveiligingsvereisten. Deze integratie valt buiten
de scope van het huidige graduaatsproject.

Daarnaast kan de compatibiliteit van de gegenereerde e-mails verder worden
onderzocht. Verdere tests in verschillende versies en platformen van Gmail,
Outlook en Apple Mail kunnen aantonen welke HTML- en CSS-elementen
betrouwbaar worden weergegeven en waar aanvullende aanpassingen nodig zijn.

Ook de editor zelf kan verder worden uitgebreid. Mogelijke uitbreidingen
zijn extra opmaakmogelijkheden, ondersteuning voor bijkomende e-mailtypes
en verbeteringen aan het beheer en de versieopvolging van templates.
Dergelijke uitbreidingen moeten worden afgewogen tegen de noodzaak om de
editor eenvoudig en beheersbaar te houden.

Ten slotte kan een toekomstig project onderzoeken hoe de PoC op een veilige
en onderhoudbare manier in de productieomgeving kan worden opgenomen.
Daarbij zijn onder meer de keuze van de definitieve opslagmethode, de
integratie met bestaande diensten, toegangsbeheer, validatie, monitoring
en het beheer van wijzigingen aan templates van belang.

De resultaten van dit graduaatsproject kunnen als uitgangspunt dienen om
de technische haalbaarheid, de resterende risico's en de vereisten voor
een mogelijke verdere ontwikkeling te beoordelen.

%---------- Bijlagen -----------------------------------------------------------

\appendix

\chapter{Projectvoorstel}

Het onderwerp van dit graduaatsproject werd vooraf vastgelegd in een
projectvoorstel, afgestemd met de mentor en goedgekeurd door de begeleider.
Dat voorstel is hieronder opgenomen.

\input{../voorstel/voorstel-inhoud}

\chapter{Repository}

De volledige broncode, de opbouw van het project en de startinstructies staan
in de repository:

\begin{center}
  \url{\repolink}
\end{center}

% Is je repository niet publiek toegankelijk (bv. omwille van een NDA of omdat
% de code eigendom blijft van het bedrijf)? Vermeld dat hier uitdrukkelijk, en
% beschrijf wat je wel kan tonen: werking, architectuur, tests, of beperkte
% fragmenten.

%%=============================================================================
%% Verantwoording van het gebruik van AI - VERPLICHT
%%=============================================================================
\chapter{\IfLanguageName{dutch}{Verantwoording van het gebruik van AI}{Use of AI}}%
\label{ch:ai}
%% wat/waarvoor/waarom gebruikt? kunnen uitleggen. 

#TODO: Vervang deze alinea door een korte inleiding: heb je AI gebruikt tijdens dit
project, en zo ja, welke rol speelde ze in je werkwijze?

\section{\IfLanguageName{dutch}{Waarvoor en hoe}{What and how}}%
\label{sec:ai-waarvoor}

%% Vul de tabel aan met je eigen gebruik. Eén regel per gebruik, niet per
%% gesprek en niet per dag.
%%
%% De twee laatste kolommen zijn de belangrijkste:
%%   - WAAR HET ZIT: wijs een plek aan in je werk - een bestand, een module, een
%%     hoofdstuk, een dia. Zat het enkel in je denkwerk en niet rechtstreeks in
%%     je resultaat, schrijf dan dat.
%%   - WAT JIJ ERMEE DEED: overgenomen, aangepast, herschreven of verworpen, en
%%     hoe je nagegaan bent dat het klopt.
%%
%% Houd dit bij TIJDENS je project, niet achteraf. In week 13 herinner je je
%% niet meer wat je in week 4 gevraagd hebt, en dat is precies te zien.

\begin{table}[htbp]
  \centering\small
  \begin{tabular}{p{0.17\textwidth}p{0.13\textwidth}p{0.2\textwidth}p{0.34\textwidth}}
    \toprule
    \textbf{Waarvoor} & \textbf{Hulpmiddel} & \textbf{Waar het in mijn werk zit} & \textbf{Wat ik ermee deed en controleerde} \\
    \midrule
    Voorbeeld: opzet van de gegevenslaag
      & naam en versie
      & \texttt{src/data/} \newline (repository)
      & Voorstel laten genereren, daarna zelf herschreven omdat de foutafhandeling ontbrak; met tests nagekeken. \\
    Voorbeeld: uitleg over een foutmelding
      & naam en versie
      & nergens rechtstreeks; hielp me het probleem begrijpen
      & Als uitleg gebruikt, de oplossing zelf geschreven na controle in de officiële documentatie. \\
    Voorbeeld: nalezen van dit rapport
      & naam en versie
      & hoofdstuk 4
      & Enkel op spelling en zinsbouw; de inhoud en de structuur zijn van mij. \\
    \bottomrule
  \end{tabular}
  \caption[Gebruik van AI in dit project]{\label{tab:ai}%
    Overzicht van het gebruik van AI in dit project.}
\end{table}

\section{\IfLanguageName{dutch}{Wat van mij is}{What is my own work}}%
\label{sec:ai-eigenwerk}

%% Twee regels die de opleiding hanteert, en die je hier bevestigt:
%%
%% 1. Code die je niet kan uitleggen, hoort niet in je oplevering. Je bent
%%    verantwoordelijk voor alles wat in je repository staat, ook voor wat een
%%    hulpmiddel geschreven heeft.
%% 2. De analyse, de keuzes, de afbakening en de tekst van dit rapport zijn van
%%    jou. AI mag je helpen formuleren; ze mag niet in jouw plaats denken.
%%
%% Geef hier ook aan hoe je in je repository zichtbaar maakt welke stukken met
%% AI tot stand kwamen - bijvoorbeeld met een korte notitie in de commit of een
%% opmerking bovenaan het bestand.

\section{\IfLanguageName{dutch}{Zorgvuldigheid en regelgeving}{Diligence and regulation}}%
\label{sec:ai-regelgeving}

%% Beantwoord hier kort de drie vragen die er in de praktijk toe doen.
%%
%% 1. HEB JE BEDRIJFSGEGEVENS IN EEN AI-HULPMIDDEL GEPLAKT?
%%    Broncode, klantgegevens, configuratie of persoonsgegevens ingeven in een
%%    publieke dienst is een doorgifte aan een derde partij. Dat kan in strijd
%%    zijn met je geheimhoudingsplicht en met de AVG (de Algemene Verordening
%%    Gegevensbescherming, ook bekend als de GDPR). Vraag je mentor wat bij
%%    jouw bedrijf toegelaten is, en noteer het antwoord hier.
%%
%% 2. ZIT ER AI IN WAT JE GEBOUWD HEBT?
%%    Dan is de Europese AI-verordening (de AI Act, verordening 2024/1689) van
%%    belang: haar verplichtingen zijn sinds augustus 2026 grotendeels van
%%    toepassing. Ga na in welke risicocategorie je toepassing valt, en of er
%%    transparantieplichten gelden - bijvoorbeeld wanneer gebruikers met een
%%    AI-systeem praten of wanneer je toepassing inhoud genereert. In beide
%%    gevallen moet de gebruiker dat weten.
%%    Gebruik je AI enkel als hulpmiddel bij het programmeren, dan valt je
%%    project daar meestal niet onder; schrijf dat dan gewoon zo op.
%%
%% 3. WAT MET AUTEURSRECHT EN LICENTIES?
%%    Gegenereerde code kan sterk lijken op bestaande code met een licentie.
%%    Vermeld hoe je daarmee omgegaan bent.

Vervang deze alinea door je antwoord op de drie vragen hierboven.

%% Heb je geen AI gebruikt? Vervang dit hele hoofdstuk dan door één zin:
%% "Bij de realisatie van dit project is geen gebruikgemaakt van generatieve AI."


%%---------- Andere bijlagen ---------------------------------------------------
% Voeg hier eventuele andere bijlagen toe.
%\input{...}

%%---------- Nawerk, bronnenlijst ----------------------------------------------

\backmatter{}

\setlength\bibitemsep{2pt}
\printbibliography[heading=bibintoc]

\end{document}
